\chapter{Langage, ensembles et preuves}
\section*{Exercice 1.1 -- Ensembles}
Dans $E=\{1,\ldots,10\}$, soit $A$ l'ensemble des multiples de $2$ et $B$ celui des multiples de $3$. Donner $A\cap B$, $A\cup B$, $\overline A$ et leurs cardinaux.

\section*{Exercice 1.2 -- Couples}
Pour $A=\{0,1\}$ et $B=\{a,b,c\}$, écrire $A\times B$. Les couples $(0,a)$ et $(a,0)$ désignent-ils la même chose ?

\section*{Exercice 1.3 -- Négations}
Nier : « $x>2$ », « $x\geq0$ et $x<5$ », « $x=-1$ ou $x=1$ ».

\section*{Exercice 1.4 -- Réciproque}
Prouver : si un entier est multiple de $6$, il est multiple de $3$. Écrire la réciproque et décider si elle est vraie.

\section*{Exercice 1.5 -- Une égalité d'ensembles}
Démontrer $A\cap(B\cup C)=(A\cap B)\cup(A\cap C)$ en raisonnant sur l'appartenance d'un élément.

\section*{Exercice 1.6 -- L'ordre des mots}
Comparer : « pour tout réel $x$, il existe un réel $y$ tel que $x+y=0$ » et « il existe un réel $y$ tel que, pour tout réel $x$, $x+y=0$ ».
